\documentclass[twoside,11pt]{article}
\usepackage{jmlr2e}

% Additional packages
\usepackage{amsmath}
\usepackage{amssymb}
\usepackage{graphicx}
\usepackage{booktabs}
\usepackage{algorithm}
\usepackage{algorithmic}
\usepackage{hyperref}
\usepackage{url}
\usepackage{xcolor}
\usepackage{listings}

% Code listings configuration
\lstset{
    language=Python,
    basicstyle=\ttfamily\small,
    keywordstyle=\color{blue},
    commentstyle=\color{green!60!black},
    stringstyle=\color{red},
    showstringspaces=false,
    breaklines=true,
    frame=single,
    numbers=left,
    numberstyle=\tiny\color{gray},
    captionpos=b
}

% JMLR heading
\jmlrheading{1}{2026}{1-30}{1/26}{0/00}{25-000}{Singh}

% Short headings for running headers
\ShortHeadings{Bayes-HDC: Probabilistic Vector Symbolic Architectures}{Singh}

\firstpageno{1}

\begin{document}

\title{Bayes-HDC: Probabilistic Vector Symbolic Architectures \\ for Hyperdimensional Computing}

% NOTE: Add co-authors once they have agreed to authorship. Listing someone
% without explicit consent is not acceptable for JMLR. The CITATION.cff at
% the repository root is the single source of truth for current authorship.
\author{\name Rajdeep Singh \email rajdeeps@usc.edu \\
       \addr Department of Computer Science\\
       University of Southern California\\
       Los Angeles, CA 90089, USA}

\maketitle

\begin{abstract}
We introduce \emph{Probabilistic Vector Symbolic Architectures} (PVSA), a framework in which every hypervector is a posterior distribution rather than a point, and every VSA primitive---binding, bundling, permutation, similarity, and retrieval---propagates the posterior's moments in closed form. Bayes-HDC is a JAX library implementing PVSA. It provides Gaussian, Dirichlet, and Mixture hypervector types, each with closed-form moment propagation under the core operations and closed-form Kullback--Leibler divergences suitable as variational objectives; a property-based treatment of the cyclic-shift group action and the diagonal vs.\ single-argument equivariance of binding; a log-space L-BFGS temperature calibrator that reduces Expected Calibration Error on standard HDC datasets; and a split-conformal classifier with the Adaptive Prediction Sets nonconformity score that returns prediction sets with marginal coverage at level $1-\alpha$ on exchangeable data. A deterministic VSA foundation---eight classical models (BSC, MAP, HRR, FHRR, BSBC, CGR, MCR, VTB), five encoders, five learning models, three memory modules, and four symbolic data structures---serves as the substrate. All types are JAX pytrees, so \texttt{jit}, \texttt{vmap}, \texttt{grad}, \texttt{pmap}, and \texttt{shard\_map} compose with every operation. The library is available under the MIT license at \url{https://github.com/rlogger/bayes-hdc}.
\end{abstract}

\begin{keywords}
Probabilistic Vector Symbolic Architectures, Bayesian Hyperdimensional Computing, Uncertainty Quantification, Calibrated Retrieval, Conformal Prediction, JAX
\end{keywords}

\section{Introduction}

Hyperdimensional computing (HDC), also known as vector symbolic architectures (VSA), is a computational framework that exploits properties of high-dimensional random vector spaces to perform symbolic reasoning through distributed representations \citep{kanerva2009hyperdimensional, gayler2003vector}. Symbols are represented as high-dimensional vectors (typically $d \geq 10{,}000$), and algebraic operations---\emph{binding}, \emph{bundling}, and \emph{permutation}---compose and manipulate these representations while preserving semantic structure \citep{plate1995holographic, kanerva2010dollar}.

Interest in HDC has grown substantially in recent years, driven by comprehensive theoretical surveys \citep{kleyko2022survey1, kleyko2022survey2}, successful applications in classification \citep{rahimi2016robust}, speech recognition \citep{imani2017voicehd}, and robotics \citep{neubert2019introduction}, and recognition of HDC's suitability for neuromorphic and edge hardware \citep{kleyko2022vector}. Neuro-symbolic architectures that combine HDC with neural networks have shown further promise \citep{hersche2023neuro}.

Despite this progress, existing HDC software treats a hypervector as a point in $\mathbb{R}^d$ with no explicit notion of confidence. Retrieval therefore returns a symbol without a calibrated probability, and classification returns a label without a coverage guarantee. In domains where HDC is being deployed---biomedical signal classification, anomaly detection, robotics---uncalibrated point predictions are a liability. Every comparable library---Torchhd \citep{heddes2023torchhd}, hdlib \citep{cumbo2023hdlib}, PyBHV \citep{vandervorst2023pybhv}---shares this limitation.

Bayes-HDC addresses the gap by lifting every hypervector to a distribution. The primary contribution is a Bayesian layer in which binding, bundling, similarity, and retrieval are defined directly on distributions, giving the user first-class control over uncertainty. The layer is built on JAX \citep{jax2018github} so that all operations are JIT-compilable, vectorisable via \texttt{vmap}, and differentiable via \texttt{grad}. Concretely, the paper contributes (i) Gaussian, Dirichlet, and Mixture hypervector types with closed-form moment propagation under bind, bundle, and permute (Section~\ref{sec:pvsa}); (ii) an explicit treatment of the cyclic-shift group action and the two distinct equivariance flavours of the HDC primitives, together with property-based verifiers; (iii) a temperature calibrator and a conformal classifier that supply calibrated probabilities and marginal-coverage guarantees on top of any HDC pipeline; and (iv) a deterministic VSA foundation of eight models, five encoders, five learning models, three memory modules, and four symbolic data structures (Section~\ref{sec:system}). All operations are pytree-native and compose with \texttt{jit} / \texttt{vmap} / \texttt{grad} / \texttt{pmap} / \texttt{shard\_map} unconditionally; Section~\ref{sec:examples} shows a single end-to-end pipeline that exercises the full stack.

\section{Probabilistic VSA}
\label{sec:pvsa}

The primary contribution is to lift the VSA algebra from $\mathbb{R}^d$ to the space $\mathcal{P}(\mathbb{R}^d)$ of distributions, in such a way that every primitive propagates the distribution's moments in closed form.

\subsection{Distribution types}

A \texttt{GaussianHV} is a frozen pytree dataclass holding a mean $\mu \in \mathbb{R}^d$ and a per-dimension variance $\sigma^2 \in \mathbb{R}_{\geq 0}^d$, modelling $X \sim \mathcal{N}(\mu, \mathrm{diag}(\sigma^2))$. The zero-variance limit recovers a deterministic hypervector exactly, which means classical pipelines lift into PVSA via \texttt{GaussianHV.from\_sample(x)} without code changes. A \texttt{DirichletHV} models a hypervector on the probability simplex through a concentration vector $\alpha \in \mathbb{R}_{>0}^d$. A \texttt{MixtureHV} composes either type into a categorical mixture, with weighted moments collapsing back to a single $\mathrm{GaussianHV}$.

\subsection{Closed-form moment propagation}

Under the standard independence assumption between operands---which classical HDC relies on for quasi-orthogonality at random initialisation---the first and second moments of the element-wise product $Z = X \cdot Y$ of two independent Gaussian hypervectors are exact:
\begin{align}
\mathbb{E}[Z] &= \mu_x \cdot \mu_y, \\
\mathrm{Var}[Z] &= \mu_x^2 \sigma_y^2 + \mu_y^2 \sigma_x^2 + \sigma_x^2 \sigma_y^2.
\end{align}
The function \texttt{bind\_gaussian} returns a \texttt{GaussianHV} carrying exactly these two moments---there is no Monte Carlo, and no delta-method approximation. The bundle case is the trivial sum-of-independent-Gaussians identity, normalised to the unit sphere via the delta method on the variance term. \texttt{permute\_gaussian} applies the cyclic-shift action to both $\mu$ and $\sigma^2$ jointly. The Dirichlet bundling identity (additive concentration) is exact and is implemented in \texttt{bundle\_dirichlet}.

\subsection{KL divergences and reparameterisation gradients}

The Kullback--Leibler divergences $D_{\mathrm{KL}}(\mathcal{N}_0 \,\|\, \mathcal{N}_1)$ and $D_{\mathrm{KL}}(\mathrm{Dir}(\alpha_0) \,\|\, \mathrm{Dir}(\alpha_1))$ are textbook closed forms; \texttt{kl\_gaussian} and \texttt{kl\_dirichlet} return them analytically. Both are differentiable end-to-end under \texttt{jax.grad}, which is what makes them usable in a variational objective. Every \texttt{GaussianHV} carries a reparameterised \texttt{sample(key)} method, so $\nabla$ composes through \texttt{bind\_gaussian}, \texttt{bundle\_gaussian}, \texttt{cleanup\_gaussian}, \texttt{inverse\_gaussian}, and the ELBO helpers in \texttt{bayes\_hdc.inference} without any user-side custom VJP.

\subsection{Group action and equivariance}
\label{sec:equivariance}

For fixed dimension $d$, cyclic shift by $k$ defines an action of $\mathbb{Z}/d$ on $\mathbb{R}^d$: $T_k(x)_i = x_{(i-k) \bmod d}$. The action is faithful, additive ($T_j \circ T_k = T_{j+k}$), and an isometry of the inner product. The HDC primitives have two distinct equivariance flavours, which the library distinguishes explicitly in \texttt{bayes\_hdc.equivariance}: \emph{diagonal} equivariance ($T_k(x \star y) = T_k(x) \star T_k(y)$ for element-wise binding) and \emph{single-argument} equivariance ($T_k(x * y) = T_k(x) * y$ for circular convolution, which acquires a $2k$ shift under the diagonal action). The module ships property-based verifiers (\texttt{verify\_shift\_equivariance}, \texttt{verify\_single\_argument\_shift\_equivariance}, \texttt{verify\_shift\_invariance}) that reject any user-defined op claiming a symmetry it does not have. This makes hypervectors a typed, symmetry-respecting substrate for downstream equivariant-architecture research.

\subsection{Uncertainty layer}

\texttt{TemperatureCalibrator} fits a single temperature $T$ to a held-out validation set by minimising the negative log-likelihood in $\log T$ via L-BFGS \citep{guo2017calibration}. The objective is convex in $\log T$ with a unique global minimum; the fitted temperature is the maximum-likelihood estimator. \texttt{ConformalClassifier} implements split-conformal prediction with the Adaptive Prediction Sets nonconformity score \citep{romano2020classification}, producing prediction sets whose marginal coverage satisfies $\Pr(y \in S(x)) \geq 1 - \alpha$ on exchangeable data, regardless of the underlying classifier. Both objects accept the public \texttt{logits} method exposed on every Bayesian classifier (\texttt{BayesianCentroidClassifier.logits}, etc.) as their canonical input.

\section{System architecture}
\label{sec:system}

PVSA is built on top of a complete deterministic VSA foundation, summarised in Table~\ref{tab:modules}. Every component is implemented from primary research papers; no code is ported from existing HDC libraries (see \texttt{ORIGINALITY.md}). All types are frozen dataclasses registered with \texttt{jax.tree\_util}, so \texttt{jit}, \texttt{vmap}, \texttt{grad}, \texttt{pmap}, and \texttt{shard\_map} compose through every operation without user-side flattening boilerplate.

\begin{table}[t]
\centering
\caption{Bayes-HDC module overview. The probabilistic layer sits in \texttt{distributions}, \texttt{bayesian\_models}, \texttt{uncertainty}, \texttt{equivariance}, \texttt{inference}, and \texttt{diagnostics}; the deterministic substrate is in the rest.}
\label{tab:modules}
\small
\begin{tabular}{lll}
\toprule
\textbf{Module} & \textbf{Purpose} & \textbf{Key components} \\
\midrule
\texttt{distributions} & PVSA primitives & GaussianHV, DirichletHV, MixtureHV; \\
 & & bind/bundle/permute/cleanup/inverse/KL \\
\texttt{bayesian\_models} & Bayesian classifiers & BayesianCentroidClassifier, BayesianAdaptiveHDC, \\
 & & StreamingBayesianHDC \\
\texttt{uncertainty} & Calibration \& coverage & TemperatureCalibrator, ConformalClassifier \\
\texttt{equivariance} & Group structure & shift, hrr\_equivariant\_bilinear, verifiers \\
\texttt{inference} & Variational helpers & elbo\_gaussian, reconstruction\_log\_likelihood\_mc \\
\texttt{diagnostics} & Posterior checks & posterior\_predictive\_check, coverage\_calibration\_check \\
\texttt{distributed} & Multi-device & pmap\_bind\_gaussian, shard\_map\_bind\_gaussian \\
\texttt{resonator} & Factorisation & probabilistic\_resonator (MCMC multi-restart) \\
\midrule
\texttt{functional} & Core HDC ops & bind, bundle, permute, similarity, n-grams, sequences \\
\texttt{vsa} & VSA models & BSC, MAP, HRR, FHRR, BSBC, CGR, MCR, VTB \\
\texttt{embeddings} & Data $\to$ HV encoders & Random, Level, Projection, Kernel (RFF), Graph \\
\texttt{structures} & Symbolic data structs & Multiset, HashTable, Sequence, Graph \\
\texttt{models} & Classical classifiers & Centroid, Adaptive, LVQ, RegLS, Clustering \\
\texttt{memory} & Associative memory & SDM, Hopfield, Attention \\
\texttt{metrics} & Capacity \& diagnostics & SNR, capacity, PR, ECE, MCE, Brier \\
\texttt{datasets} & Standard benchmarks & 11 loaders (iris, MNIST, UCIHAR, EMG, \ldots) \\
\bottomrule
\end{tabular}
\end{table}

\subsection{Comparison with existing software}

Table~\ref{tab:comparison} contrasts Bayes-HDC's feature surface with the three other actively maintained HDC libraries. The framing axes are derived from the underlying framework choice (JAX vs.\ PyTorch) and the presence or absence of a probabilistic layer; they are descriptive of capability, not a head-to-head performance comparison.

\begin{table}[t]
\centering
\caption{Feature comparison of HDC/VSA software libraries.}
\label{tab:comparison}
\small
\begin{tabular}{lcccccc}
\toprule
\textbf{Library} & \textbf{VSA models} & \textbf{XLA} & \textbf{TPU} & \textbf{PVSA} & \textbf{Conformal} & \textbf{Calibration} \\
\midrule
Torchhd \citep{heddes2023torchhd} & 10+ & $\times$ & $\times$ & $\times$ & $\times$ & $\times$ \\
hdlib \citep{cumbo2023hdlib} & 2 & $\times$ & $\times$ & $\times$ & $\times$ & $\times$ \\
PyBHV \citep{vandervorst2023pybhv} & 1 & $\times$ & $\times$ & $\times$ & $\times$ & $\times$ \\
\textbf{Bayes-HDC} & 8 & \checkmark & \checkmark & \checkmark & \checkmark & \checkmark \\
\bottomrule
\end{tabular}
\end{table}

\section{Usage}
\label{sec:examples}

A single end-to-end pipeline exercises the PVSA layer: build a probabilistic centroid classifier, fit a temperature on held-out logits, wrap the calibrated probabilities in a conformal predictor, and read off coverage-guaranteed prediction sets. The same pattern drives the application examples shipped with the library (gesture recognition, activity recognition, language identification).

\begin{lstlisting}[caption={A PVSA classification pipeline with calibration and coverage guarantees.},label={lst:pvsa}]
from bayes_hdc import (
    BayesianCentroidClassifier, ConformalClassifier,
    RandomEncoder, MAP, TemperatureCalibrator,
)

vsa = MAP.create(dimensions=10_000)
encoder = RandomEncoder.create(
    num_features=F, num_values=L, dimensions=10_000,
    vsa_model=vsa, key=key,
)
hv_tr = encoder.encode_batch(X_tr_idx)
hv_ca = encoder.encode_batch(X_ca_idx)
hv_te = encoder.encode_batch(X_te_idx)

clf = BayesianCentroidClassifier.create(
    num_classes=K, dimensions=10_000,
).fit(hv_tr, y_tr)

# `logits` is the canonical pre-softmax score; it is the public input
# to TemperatureCalibrator.fit and ConformalClassifier.fit.
calib = TemperatureCalibrator.create().fit(clf.logits(hv_ca), y_ca)
probs_ca = calib.calibrate(clf.logits(hv_ca))
probs_te = calib.calibrate(clf.logits(hv_te))

conf = ConformalClassifier.create(alpha=0.1).fit(probs_ca, y_ca)
sets = conf.predict_set(probs_te)         # (n, K) bool mask
cov  = conf.coverage(probs_te, y_te)      # >= 0.9 by construction
\end{lstlisting}

The classifier object is itself a JAX pytree: all four library calls---\texttt{fit}, \texttt{logits}, \texttt{calibrate}, \texttt{predict\_set}---compile under \texttt{jax.jit} without user-side flattening. \texttt{BayesianCentroidClassifier.var} (per-class posterior variance) supplies an out-of-distribution signal via posterior Mahalanobis distance, which is unavailable from a deterministic centroid classifier. \texttt{StreamingBayesianHDC.fit} runs as a single \texttt{jax.lax.scan} so the same end-to-end pipeline works on a non-stationary stream with bounded $O(K \cdot d)$ memory.

\section{Related work}
\label{sec:related}

The theoretical foundations of HDC/VSA were established by \citet{kanerva1988sparse,kanerva2009hyperdimensional} and \citet{plate1995holographic}; recent surveys \citep{kleyko2022survey1,kleyko2022survey2} cover models, data transformations, and applications. The most mature HDC software libraries are Torchhd \citep{heddes2023torchhd} (PyTorch, broad VSA coverage), hdlib \citep{cumbo2023hdlib} (domain-oriented bioinformatics), and PyBHV \citep{vandervorst2023pybhv} (Boolean hypervectors). None ships a probabilistic layer, a temperature calibrator, or a conformal predictor; Bayes-HDC fills that gap on top of a JAX substrate. Calibration via post-hoc temperature scaling follows \citet{guo2017calibration}; conformal prediction with Adaptive Prediction Sets follows \citet{romano2020classification}. HDC's group-theoretic structure connects to current work on equivariant neural functionals and weight-space learning, where the cyclic-shift action and circular-convolution binding provide a natural bilinear primitive.

\section{Conclusion}
\label{sec:conclusion}

Bayes-HDC is a JAX library for hyperdimensional computing that lifts the standard VSA algebra to a probabilistic algebra over distributions. The closed-form moment propagation under binding and bundling, the property-based equivariance verifiers for the cyclic-shift action, the post-hoc temperature calibrator, and the split-conformal classifier with marginal-coverage guarantees together form a small, typed surface on which calibrated and coverage-guaranteed retrieval are first-class operations rather than after-the-fact wrappers. The deterministic VSA foundation is implemented from primary research and supports the library as a standalone classical-HDC toolkit. Planned work includes variational codebook learning via the reparameterised ELBO, sparse PVSA posteriors, and integration with \texttt{flax.nnx} for use as Bayesian HDC modules inside larger neural pipelines.

\section*{Acknowledgments}

We thank the authors of the foundational HDC research papers whose work this project builds on---P. Kanerva, T.~A. Plate, R.~W. Gayler, D. Kleyko, A. Rahimi, M. Imani, and F.~T. Sommer among many others---and the JAX development team at Google for the underlying numerical-computing substrate. Torchhd, hdlib, and PyBHV are cited for head-to-head empirical comparison and for their prior framing of the HDC software-engineering problem; no code in this library is derived from them.

\bibliographystyle{plainnat}
\begin{thebibliography}{99}

\bibitem[Frady et~al.(2020)]{frady2020resonator}
E.~P. Frady, S.~J. Kent, B.~A. Olshausen, and F.~T. Sommer.
\newblock Resonator networks, 1: An efficient solution for factoring high-dimensional, distributed representations of data structures.
\newblock \textit{Neural Computation}, 32(12):2311--2331, 2020.

\bibitem[Gayler(2003)]{gayler2003vector}
R.~W. Gayler.
\newblock Vector symbolic architectures answer {J}ackendoff's challenges for cognitive neuroscience.
\newblock In \textit{Joint International Conference on Cognitive Science}, pages 133--138, 2003.

\bibitem[Heddes et~al.(2023)]{heddes2023torchhd}
M.~Heddes, I.~Nunes, P.~Verg\'es, D.~Kleyko, D.~Abraham, T.~Givargis, A.~Nicolau, and A.~Veidenbaum.
\newblock Torchhd: An open source {P}ython library to support research on hyperdimensional computing and vector symbolic architectures.
\newblock \textit{Journal of Machine Learning Research}, 24(255):1--10, 2023.

\bibitem[Hersche et~al.(2023)]{hersche2023neuro}
M.~Hersche, M.~Zeqiri, L.~Benini, A.~Sebastian, and A.~Rahimi.
\newblock A neuro-vector-symbolic architecture for solving {R}aven's progressive matrices.
\newblock \textit{Nature Machine Intelligence}, 5(4):363--375, 2023.

\bibitem[Imani et~al.(2017)]{imani2017voicehd}
M.~Imani, D.~Kong, A.~Rahimi, and T.~Rosing.
\newblock {VoiceHD}: Hyperdimensional computing for efficient speech recognition.
\newblock In \textit{International Conference on Rebooting Computing (ICRC)}, pages 1--8, 2017.

\bibitem[Imani et~al.(2021)]{imani2021onlinehd}
M.~Imani, Y.~Kim, S.~Patil, and T.~Rosing.
\newblock {OnlineHD}: Robust, efficient, and single-pass online learning using hyperdimensional system.
\newblock In \textit{Design, Automation \& Test in Europe (DATE)}, pages 56--61, 2021.

\bibitem[JAX Team(2018)]{jax2018github}
JAX Team.
\newblock {JAX}: Composable transformations of {P}ython+{N}um{P}y programs.
\newblock \url{https://github.com/google/jax}, 2018.

\bibitem[Kanerva(1988)]{kanerva1988sparse}
P.~Kanerva.
\newblock \textit{Sparse Distributed Memory}.
\newblock MIT Press, Cambridge, MA, 1988.

\bibitem[Kanerva(1997)]{kanerva1997fully}
P.~Kanerva.
\newblock Fully distributed representation.
\newblock In \textit{Real World Computing Symposium (RWC)}, pages 358--365, 1997.

\bibitem[Kanerva(2009)]{kanerva2009hyperdimensional}
P.~Kanerva.
\newblock Hyperdimensional computing: An introduction to computing in distributed representation with high-dimensional random vectors.
\newblock \textit{Cognitive Computation}, 1(2):139--159, 2009.

\bibitem[Kanerva(2010)]{kanerva2010dollar}
P.~Kanerva.
\newblock What's the dollar of {M}exico? {H}yperdimensional computing in question-answering.
\newblock Technical report, Redwood Center for Theoretical Neuroscience, 2010.

\bibitem[Kim et~al.(2024)]{kim2024flash}
Y.~Kim, M.~Imani, B.~Moons, and T.~Rosing.
\newblock {FLASH}: Hyperdimensional computing with holographic and adaptive encoder.
\newblock \textit{Frontiers in Neuroscience}, 18:1--15, 2024.

\bibitem[Kleyko et~al.(2022a)]{kleyko2022survey1}
D.~Kleyko, D.~A. Rachkovskij, E.~Osipov, and A.~Rahimi.
\newblock A survey on hyperdimensional computing aka vector symbolic architectures, {P}art {I}: Models and data transformations.
\newblock \textit{ACM Computing Surveys}, 55(6):1--40, 2022.

\bibitem[Kleyko et~al.(2022b)]{kleyko2022survey2}
D.~Kleyko, D.~A. Rachkovskij, E.~Osipov, and A.~Rahimi.
\newblock A survey on hyperdimensional computing aka vector symbolic architectures, {P}art {II}: Applications, cognitive models, and challenges.
\newblock \textit{ACM Computing Surveys}, 55(9):1--52, 2022.

\bibitem[Kleyko et~al.(2022c)]{kleyko2022vector}
D.~Kleyko, M.~Davies, E.~P. Frady, P.~Kanerva, S.~J. Kent, B.~A. Olshausen, E.~Osipov, J.~M. Rabaey, D.~A. Rachkovskij, A.~Rahimi, and F.~T. Sommer.
\newblock Vector symbolic architectures as a computing framework for emerging hardware.
\newblock \textit{Proceedings of the IEEE}, 110(10):1538--1571, 2022.

\bibitem[Cumbo et~al.(2023)]{cumbo2023hdlib}
F.~Cumbo, E.~Weitschek, and D.~Blankenberg.
\newblock hdlib: A {P}ython library for designing vector-symbolic architectures.
\newblock \textit{Journal of Open Source Software}, 8(89):5704, 2023.

\bibitem[Neubert et~al.(2019)]{neubert2019introduction}
P.~Neubert, S.~Schubert, and P.~Protzel.
\newblock An introduction to hyperdimensional computing for robotics.
\newblock \textit{KI -- K{\"u}nstliche Intelligenz}, 33(4):319--330, 2019.

\bibitem[Plate(1991)]{plate1991holographic}
T.~A. Plate.
\newblock Holographic reduced representations: Convolution algebra for compositional distributed representations.
\newblock In \textit{International Joint Conference on Artificial Intelligence (IJCAI)}, pages 30--35, 1991.

\bibitem[Plate(1995)]{plate1995holographic}
T.~A. Plate.
\newblock Holographic reduced representations.
\newblock \textit{IEEE Transactions on Neural Networks}, 6(3):623--641, 1995.

\bibitem[Rahimi and Recht(2007)]{rahimi2007random}
A.~Rahimi and B.~Recht.
\newblock Random features for large-scale kernel machines.
\newblock In \textit{Advances in Neural Information Processing Systems (NIPS)}, volume~20, pages 1--8, 2007.

\bibitem[Rahimi et~al.(2016)]{rahimi2016robust}
A.~Rahimi, P.~Kanerva, and J.~M. Rabaey.
\newblock A robust and energy-efficient classifier using brain-inspired hyperdimensional computing.
\newblock In \textit{International Symposium on Low Power Electronics and Design (ISLPED)}, pages 64--69, 2016.

\bibitem[Vandervorst(2023)]{vandervorst2023pybhv}
A.~Vandervorst.
\newblock {PyBHV}: Boolean hypervectors for experiments in hyperdimensional computing.
\newblock \url{https://github.com/Adam-Vandervorst/PyBHV}, 2023.

\bibitem[Ramsauer et~al.(2020)]{ramsauer2020hopfield}
H.~Ramsauer, B.~Sch{\"a}fl, J.~Lehner, et~al.
\newblock {H}opfield networks is all you need.
\newblock arXiv preprint arXiv:2008.02217, 2020.

\end{thebibliography}

% Implementation notes (JAX dataclass registration, version-compatibility
% layer, ruff/mypy/CI matrix) are documented in DESIGN.md at the
% repository root rather than reproduced here, to respect the MLOSS
% four-page body limit.

\end{document}
